\session{13}{11/14}{2限}{4}{AIのリスクとガバナンス}

\goals{
  \item AIのリスク（バイアス、プライバシー、著作権、情報セキュリティ）を具体例とともに説明できる。
  \item AI規制の動向（EU AI Act、日本の枠組み）の基本的な考え方を理解する。
  \item 企業のAIガバナンス体制に何が必要かを説明し、方針案を書ける。
}

\guidesec{解説}

\topic{AIのリスク}
AIの活用には、他の情報技術と共通するリスクと、AI固有のリスクがあります。主なものを整理します。
\par\noindent\begin{gtabh}{colspec={Q[l,wd=26mm] X[l] X[l]}}
  リスク & 内容 & 例 \\
  バイアスと公平性 & 学習データの偏りがAIの判断に引き継がれ、特定の属性の人に不利な結果を生む & 過去の採用データで学習した選考モデルが女性を低く評価する \\
  プライバシー & 個人情報の不適切な収集・利用、入力したデータの外部流出、推論による個人の特定 & 顧客の相談内容をAIに入力し、学習に使われる \\
  著作権 & 学習に使った著作物の扱い、出力が既存の著作物に類似する問題 & 生成した画像が既存のキャラクターに酷似する \\
  情報セキュリティ & 機密情報の流出、無許可のAI利用（シャドーAI）、プロンプトインジェクション & 従業員が個人のアカウントで社外秘の資料を要約させる \\
  誤情報と品質 & ハルシネーションによる誤った情報の提供、責任の所在の不明確さ & チャットボットが誤った返金条件を案内する \\
  依存と技能の低下 & AIへの過度な依存で、人の判断力や技能が維持されない & 若手が自分で考えずにAIの案を使い続ける \\
\end{gtabh}

\textbf{バイアス}の根本原因は、学習データの偏りです。過去のデータは、過去の社会や組織の偏りを含んでいます。
AIはそれを「正解」として学習するため、偏りを再生産し、しかも「客観的な判断」という外見で正当化してしまいます。
対策は、学習データの点検、属性ごとの結果の比較（監査）、影響の大きい判断での人の関与（第11回）です。

\textbf{著作権}については、学習段階と利用段階を分けて考えます。日本の著作権法では、情報解析のための著作物の利用は一定の条件で認められていますが、
出力が既存の著作物に類似し、それに依拠していると判断されれば、利用段階で侵害になる可能性があります。
現行法では AI は著作者になれず、AI の出力に著作権が発生するかは人の創作的関与によります。生成した画像や文章を公開するときは、既存の作品に似ていないかを確認します。

\textbf{シャドーAI}は、組織が把握・許可していない AI ツールを従業員が業務に使うことで、シャドーIT の派生概念です。
禁止するだけでは地下化するため、安全な公式ツールを提供し、使ってよい範囲を示し、教育することが対策になります。
\textbf{プロンプトインジェクション}は、AIに読み込ませる文書やWebページに不正な指示を仕込み、AIの動作を乗っ取る攻撃です。
AIエージェント（第5回）が外部の情報を読んでツールを操作する場合、特に注意が必要です。

\topic{AI規制の動向}
AIに関する規制は、各国・地域で整備が進んでいます。基本的な考え方を押さえておきます。
\begin{itemize}
  \item \textbf{EU AI Act}　2024年に成立した包括的なAI規制法で、AIの用途をリスクの高さで「許容できないリスク（禁止）」「高リスク（厳格な義務）」
        「限定的リスク（透明性の義務）」「最小リスク（義務なし）」に分類し、リスクに応じた義務を課すリスクベースの規制です。
        義務は2025年から段階的に適用されています。EU 域内に製品やサービスを提供する日本企業にも適用されます。
  \item \textbf{日本}　法的な拘束力を持つ包括的な規制ではなく、政府の「AI事業者ガイドライン」などによる、事業者の自主的な取り組みを促す枠組みが基本です。
        2025年にはAIの研究開発と活用を推進するための法律（AI推進法）が成立し、政府の戦略本部と基本計画の枠組みが整えられました。
        個人情報保護法や著作権法など既存の法律は、AIの利用にもそのまま適用されます。
  \item \textbf{国際的な動き}　OECD の AI 原則、G7 の広島 AI プロセスなど、国際的な原則づくりも進んでいます。
\end{itemize}
規制の内容は変わり続けるため、最新の情報は官公庁の公式サイトで確認してください。経営の視点で重要なのは、
「法律で禁止されていないから使ってよい」ではなく、自社がどんなリスクをどこまで取るかを自ら決めて説明できることです。

\topic{企業のAIガバナンス}
AIガバナンスは、AIを適切に活用するために組織が整える方針・体制・プロセスです。主な要素を挙げます。
\begin{itemize}
  \item \textbf{利用方針}　使ってよい用途と使ってはいけない用途、入力してよい情報の範囲、出力の確認の義務、AI利用の明示。
  \item \textbf{リスク評価}　AIの用途ごとに、影響の大きさと発生の可能性を評価し、人の関与の度合い（第11回）を決める。
  \item \textbf{責任者と体制}　AI活用の責任者、推進と監視の役割分担、相談窓口。法務・情報セキュリティ・現場の連携。
  \item \textbf{教育}　全従業員への基礎教育と、利用者への実践教育。ハルシネーション、情報の扱い、著作権。
  \item \textbf{監査と見直し}　利用状況の記録、出力の品質やバイアスの定期的な点検、事故が起きたときの報告と対応、方針の更新。
  \item \textbf{取引先・顧客への説明}　AIを使っていることの開示、顧客からの問い合わせへの対応。
\end{itemize}
本講義の AI Policy（使う・検証する・区別する）とAI活用記録は、個人の水準でのガバナンスの実践だと言えます。

\topic{キーワード}
\par\noindent\begin{keywords}{}
  バイアス & 学習データの偏りがAIの判断に引き継がれ、特定の属性に不利な結果を生むこと。 \\
  シャドーAI & 組織の許可なく従業員がAIツールを業務に使うこと。情報漏えいのリスク。 \\
  プロンプトインジェクション & AIが読む文書に不正な指示を仕込み、動作を乗っ取る攻撃。 \\
  EU AI Act & リスクの高さに応じて義務を変えるリスクベースの包括的AI規制。 \\
  AI事業者ガイドライン & 日本政府が示す、AI開発者・提供者・利用者向けの自主的な取り組みの指針。 \\
  AIガバナンス & 利用方針、リスク評価、責任者・体制、教育、監査を整え、AIを適切に活用するしくみ。 \\
\end{keywords}

\guidesec{演習課題}

\exercise{1}{AI演習：AIの出力に含まれる偏りや問題点を実際に検出し、対策を考える}
\instr{次のプロンプトを Copilot に入力し、出力に含まれる偏り（性別、年齢、国籍、職業などに関する固定観念）、根拠のない断定、著作権上の懸念を探してください。見つけた問題ごとに、「なぜ生じたと考えられるか」と「プロンプトや確認の手順でどう防ぐか」を書きます。同じプロンプトを2〜3回実行して、出力の違いも観察してください。}
\begin{promptbox}[偏りを観察する]
新しく開く家電量販店の店長候補として、次の4人のうち誰が最も適任かを順位付けし、理由を述べてください。A：45歳男性、店舗勤務20年。B：32歳女性、店舗勤務8年、育児中。C：28歳男性、EC 部門勤務5年。D：50歳女性、他業種の店長経験10年。
\end{promptbox}
\instr{次に、「それぞれの人物について、店長に必要な能力の観点から、順位を付けずに強みと確認すべき点を挙げてください」と聞き直し、出力がどう変わるかを比べてください。}

\exercise{2}{リスクの一覧表を作る}
\instr{「顧客からの問い合わせメールに対して、生成AIが回答案を作り、担当者が確認して送信する」という業務について、解説の6つのリスクそれぞれが具体的にどう現れるかと、対策を表にしてください。}

\exercise{3}{AI利用方針の案を書く}
\instr{従業員100人の会計事務所を想定し、従業員向けのAI利用方針を「目的」「使ってよい用途」「使ってはいけない用途」「入力してよい情報・いけない情報」「出力の確認と明示」「相談窓口と違反時の扱い」の6項目、各2〜3行で書いてください。}

\guidesec{模範解答}

\begin{answer}{演習1}
観察の例：
\begin{itemize}
  \item 偏り：A（男性・勤務20年）を1位にし、B について「育児中であるため勤務時間に制約がある可能性」と、本人に確認していない事情を不利な要素として挙げた。
        D について「他業種のため家電の知識が不足」と断定した。実行のたびに順位は変わるが、B が1位になることは少なかった。
  \item 根拠のない断定：「店長には長時間の勤務が必要」など、店舗の条件を知らないのに前提を置いている。
  \item 原因：学習データに含まれる社会の固定観念（育児と管理職、年齢と適応力）が、そのまま判断の根拠として現れた。また、プロンプトが「順位付け」を求めたため、
        AIは与えられた属性だけで差をつけざるを得なかった。
  \item 対策：人事に関わる判断を AI に順位付けさせない（Human-in-the-loop、第11回）。属性ではなく職務に必要な能力と実績で評価する枠組みを先に決め、
        AIには「確認すべき点の整理」だけを求める。出力を属性ごとに点検する手順を置く。
\end{itemize}
聞き直した結果：順位がなくなり、各人の強み（A：店舗運営の経験、B：現場と最新の店舗事情、C：EC との連携、D：他業種の視点）と
確認すべき点（本人の希望、勤務条件、マネジメント経験の内容）が並列に示された。問いの立て方で、偏りの現れ方が大きく変わることが分かる。
\end{answer}

\begin{answer}{演習2}
\par\noindent\begin{gtabh}{colspec={Q[l,wd=26mm] X[l] X[l]}}
  リスク & この業務での現れ方 & 対策 \\
  バイアス & 顧客の名前や文体から属性を推定し、対応の丁寧さや提案内容に差が出る & 回答の基準（返金条件、対応の水準）を規程として与え、属性に依存しない回答にする。属性別に回答の品質を点検する。 \\
  プライバシー & 顧客のメールに含まれる個人情報がAIに送られ、学習に使われる & 学習に使わない法人契約。氏名・連絡先をマスクしてから入力する仕組み。 \\
  著作権 & 回答案に他社の説明文やマニュアルの文章がそのまま含まれる & 自社の規程・FAQ を根拠に回答させ（RAG）、外部文書の引用を禁じる。 \\
  情報セキュリティ & 担当者が公式ツール以外を使う。顧客のメールに仕込まれた指示でAIが誤った回答を出す（プロンプトインジェクション） & 公式ツールの提供と利用の記録。メールの内容を「データ」として扱い、指示として実行しない設計。 \\
  誤情報と品質 & 返金条件や納期を誤って案内し、顧客との紛争になる & 事実に関わる部分（金額、期日、条件）は担当者が規程と照合する確認項目を設ける。 \\
  依存と技能の低下 & 担当者が回答案を読まずに送信するようになる & 確認した項目のチェックを記録する。定期的に、AIなしで回答を書く訓練を行う。 \\
\end{gtabh}
\end{answer}

\begin{answer}{演習3}
\begin{itemize}
  \item 目的：顧問先への対応の質と速さを高め、定型業務の時間を専門的な助言に振り向けるために、生成AIを安全に活用する。
  \item 使ってよい用途：法令・通達の要約と検索の補助、社内文書・メールの下書き、会議の議事録作成、表計算の数式やコードの作成。
  \item 使ってはいけない用途：顧問先への正式な回答・申告書類をAIの出力のまま提出すること。顧問先や職員の評価・選別をAIに判断させること。
  \item 入力してよい情報・いけない情報：公開情報と、自分が作成した文書は入力してよい。顧問先の氏名・財務データ・申告内容など識別可能な情報は、事務所が契約した公式ツール以外には入力しない。公式ツールでも、必要最小限に絞る。
  \item 出力の確認と明示：法令の条文番号、数値、期日は必ず原典で確認する。顧問先に提出する文書にAIを使った場合は、内部の作業記録に使用箇所を記録する。
  \item 相談窓口と違反時の扱い：判断に迷う場合は所長補佐に相談する。誤って機密情報を入力した場合は、速やかに報告することを優先し、報告者を不利に扱わない。方針は半年ごとに見直す。
\end{itemize}
\end{answer}

\guidesec{出席確認クイズの解説}
講義サイトの出席確認ページ（第13回）の5問です。
% 出席確認クイズ 第13回　AIのリスクとガバナンス
%   attendance/attendance13.html から生成（解説は本ガイド用に加筆）。

\quizq{1}{AIの「バイアス」の主な原因として適切なものはどれでしょうか?}%
  {電力不足}%
  {画面の大きさ}%
  {\correct{学習データの偏り}}%
  {通信速度}%
  {正解は (c)。学習データに偏りがあると、AIの判断もその偏りを引き継ぎます。たとえば過去の採用データに性別の偏りがあれば、そこから学習した選考モデルも同じ偏りを持ちます。(a)(b)(d) はバイアスの原因ではありません。}

\quizq{2}{「シャドーAI」の説明として適切なものはどれでしょうか?}%
  {公式に承認されたAI}%
  {\correct{組織の許可なく従業員がAIツールを業務に使うこと}}%
  {夜間に稼働するAI}%
  {画像の影を処理するAI}%
  {正解は (b)。シャドーAIは、組織が把握・許可していない AI ツールを従業員が業務に使うことで、機密情報の漏えいなどのリスクがあります。シャドーIT（無許可の IT 利用）の派生概念です。禁止するだけでは地下化するため、安全な公式ツールの提供と教育が対策になります。}

\quizq{3}{EUのAI規制(AI Act)の基本的な考え方として適切なものはどれでしょうか?}%
  {規制は一切しない}%
  {企業の規模だけで決める}%
  {すべてのAIを禁止する}%
  {\correct{リスクの高さに応じて義務や規制の強さを変える}}%
  {正解は (d)。EU AI Act は、AIの用途を「許容できないリスク」「高リスク」「限定的リスク」「最小リスク」に分類し、リスクの高さに応じた義務を課すリスクベースの規制です。(a)(b)(c) はいずれも実際の枠組みとは異なります。}

\quizq{4}{生成AIと著作権に関する説明として適切なものはどれでしょうか?}%
  {\correct{出力が既存の著作物に類似していると、著作権侵害になる可能性がある}}%
  {著作権はAIには関係ない}%
  {著作権はAIが持つ}%
  {AIが作ったものはすべて自由に使える}%
  {正解は (a)。生成AIの出力であっても、既存の著作物と類似し、それに依拠していると判断されれば著作権侵害になる可能性があります。(b)(d) は誤りで、(c) 現行法では AI は著作者になれません。学習段階と利用（出力）段階で考え方が異なる点にも注意が必要です。}

\quizq{5}{企業のAIガバナンスに含まれるものとして適切なものはどれでしょうか?}%
  {ルールを作らない}%
  {AIの利用を隠す}%
  {全社員に無制限の利用を許可する}%
  {\correct{利用方針の策定、リスク評価、責任者・体制の明確化}}%
  {正解は (d)。AIガバナンスは、利用方針の策定、リスク評価、責任者・体制の明確化、教育、監査などを整え、AIを適切に活用するためのしくみです。(a)(c) はガバナンスの不在、(b) は透明性の否定です。}

